\documentclass[10pt,a4paper]{amsart}
\usepackage{amsmath,amssymb,amsthm,graphicx}
% Retained preamble for the Stacks Project chapter morphisms.tex (Morphisms of Schemes).
% The chapter file pinned-original-source.tex carries only \input{preamble} on its first line, so no
% preamble text lives inside the chapter itself.  The Stacks Project uses one single shared preamble for
% every chapter; the snapshot stages it verbatim in sources/stacks-chow/preamble.tex.  The block below is
% copied mechanically, line for line, from that file's lines 14-18 (the Xy-pic setup) and lines 154-217
% (the native theorem declarations together with the mathematical macros).  Nothing was added, removed or
% redefined.  The wrapper class and the amsmath/amssymb/amsthm/graphicx base packages are supplied by the
% assembler, so the class/\externaldocument part of the shared preamble is not repeated here.
% For commutative diagrams we use Xy-pic
\usepackage[all]{xy}

% We use 2cell for 2-commutative diagrams.
\xyoption{2cell}
% Theorem environments.
%
\theoremstyle{plain}
\newtheorem{theorem}[subsection]{Theorem}
\newtheorem{proposition}[subsection]{Proposition}
\newtheorem{lemma}[subsection]{Lemma}

\theoremstyle{definition}
\newtheorem{definition}[subsection]{Definition}
\newtheorem{example}[subsection]{Example}
\newtheorem{exercise}[subsection]{Exercise}
\newtheorem{situation}[subsection]{Situation}

\theoremstyle{remark}
\newtheorem{remark}[subsection]{Remark}
\newtheorem{remarks}[subsection]{Remarks}

\numberwithin{equation}{subsection}

% Macros
%
\def\lim{\mathop{\mathrm{lim}}\nolimits}
\def\colim{\mathop{\mathrm{colim}}\nolimits}
\def\Spec{\mathop{\mathrm{Spec}}}
\def\Hom{\mathop{\mathrm{Hom}}\nolimits}
\def\Ext{\mathop{\mathrm{Ext}}\nolimits}
\def\SheafHom{\mathop{\mathcal{H}\!\mathit{om}}\nolimits}
\def\SheafExt{\mathop{\mathcal{E}\!\mathit{xt}}\nolimits}
\def\Sch{\mathit{Sch}}
\def\Mor{\mathop{\mathrm{Mor}}\nolimits}
\def\Ob{\mathop{\mathrm{Ob}}\nolimits}
\def\Sh{\mathop{\mathit{Sh}}\nolimits}
\def\NL{\mathop{N\!L}\nolimits}
\def\CH{\mathop{\mathrm{CH}}\nolimits}
\def\proetale{{pro\text{-}\acute{e}tale}}
\def\etale{{\acute{e}tale}}
\def\QCoh{\mathit{QCoh}}
\def\Ker{\mathop{\mathrm{Ker}}}
\def\Im{\mathop{\mathrm{Im}}}
\def\Coker{\mathop{\mathrm{Coker}}}
\def\Coim{\mathop{\mathrm{Coim}}}

% Boxtimes
%
\DeclareMathSymbol{\boxtimes}{\mathbin}{AMSa}{"02}

%
% Macros for moduli stacks/spaces
%
\def\QCohstack{\mathcal{QC}\!\mathit{oh}}
\def\Cohstack{\mathcal{C}\!\mathit{oh}}
\def\Spacesstack{\mathcal{S}\!\mathit{paces}}
\def\Quotfunctor{\mathrm{Quot}}
\def\Hilbfunctor{\mathrm{Hilb}}
\def\Curvesstack{\mathcal{C}\!\mathit{urves}}
\def\Polarizedstack{\mathcal{P}\!\mathit{olarized}}
\def\Complexesstack{\mathcal{C}\!\mathit{omplexes}}
% \Pic is the operator that assigns to X its picard group, usage \Pic(X)
% \Picardstack_{X/B} denotes the Picard stack of X over B
% \Picardfunctor_{X/B} denotes the Picard functor of X over B
\def\Pic{\mathop{\mathrm{Pic}}\nolimits}
\def\Picardstack{\mathcal{P}\!\mathit{ic}}
\def\Picardfunctor{\mathrm{Pic}}
\def\Deformationcategory{\mathcal{D}\!\mathit{ef}}

% --- appended 2026-10-12 (Chief): the project own definitions of the reference/slogan/history
% display environments, copied mechanically from the shared Stacks preamble lines 8-12
% (\usepackage{verbatim} plus the three \newenvironment lines, which the chapter preamble keeps
% verbatim from the project preamble). Nothing added, removed or redefined: these environments
% carry no mathematical content and the project itself typesets them as comments. Without these
% three lines a counted statement whose own text uses them fails with Environment slogan undefined.
\usepackage{verbatim}
\newenvironment{reference}{\comment}{\endcomment}
\newenvironment{slogan}{\comment}{\endcomment}
\newenvironment{history}{\comment}{\endcomment}
\usepackage{cleveref}
\crefname{originallemma}{original lemma}{original lemmas}
\Crefname{originallemma}{Original lemma}{Original lemmas}
\makeatletter
\newlabel{morphisms-lemma-birational-birational}{{0BAC}{source}{}{sourcepointer-1}{}}
\newlabel{morphisms-lemma-birational-birational@cref}{{[originallemma][0][]0BAC}{[1][0][]source}{}{}{}}
\makeatother
\crefname{originallemma}{original lemma}{original lemmas}
\Crefname{originallemma}{Original lemma}{Original lemmas}
\makeatletter
\newlabel{lemma-birational-generic-fibres}{{0BAB}{source}{}{sourcepointer-2}{}}
\newlabel{lemma-birational-generic-fibres@cref}{{[originallemma][0][]0BAB}{[1][0][]source}{}{}{}}
\makeatother
\title{A composition of birational morphisms is birational}
\author{The Stacks Project Authors}
\date{stacks-project:morphisms.tex:local-snapshot-20261012; GNU FDL 1.2 or later}
\begin{document}
\maketitle
\setcounter{subsection}{1}
\noindent\textit{Editorial note.} All delivered mathematical text is the author's own, copied from the pinned Stacks Project chapter file morphisms.tex; no statement or proof step was rewritten, reordered or removed.\par\medskip
\noindent\textit{Editorial note.} The chapter file pinned-original-source.tex carries only \textbackslash{}input\{preamble\} on its first line, so no preamble text lives inside the chapter. The retained file sources/stacks-morphisms/retained-preamble.tex is a mechanical, line-for-line copy of the project's single shared preamble as staged in the snapshot. Nothing was added, removed or redefined, and the excerpt serves only as preamble, never as a span coordinate.\par\medskip
\noindent\textit{Editorial note.} Both counted spans start at the \textbackslash{}begin line of their environment, so the statement and the author's proof are delivered whole; surrounding prose of the section is deliberately outside the counted spans and is carried as standard prerequisites instead.\par\medskip
\noindent\textit{Editorial note.} The citation fields (source\_id, source\_version, authors, source\_url, license.name, source\_title) are identical to those of every other card this lane has produced from this chapter, so the batch selector treats them as one source.\par\medskip
\noindent\textit{Editorial note.} Read in full under the Lead's 12:4x rule (screening hit re-read, pb 1475B). The author writes verbatim 'There is an immediate reduction to the case where X and Y are irreducible which we omit.' - the reduction is to the standard irreducible case and the rest of the proof is carried out on the minimal primes; recorded here as the author's own wording.\par\medskip
\noindent\textit{Allowed standard prerequisites (editorial).}\begin{itemize}
\item Birationality described through the unique minimal primes and the induced extension of function fields. Lemma \textbackslash{}ref\{lemma-birational-integral\} and the surrounding discussion, used at proof lines 13087-13102. Gives the criterion applied to both maps and to their composition.
\end{itemize}\medskip
\ifdefined\hypertarget\hypertarget{sourcepointer-1}{}\fi
\noindent\textit{Source reference.} Original lemma 0BAC: focused-seeds/AGP/sources/stacks-snapshot-20261011/tags.txt, source lines 14205--14205. The lemma's own canonical tag, quoted so that readers can resolve the label inside the delivered text; it points at the counted statement itself and carries no proof content.\par\medskip
\ifdefined\hypertarget\hypertarget{sourcepointer-2}{}\fi
\noindent\textit{Source reference.} Original lemma 0BAB: focused-seeds/AGP/sources/stacks-snapshot-20261011/tags.txt, source lines 14204--14204. Navigation pointer to a statement the counted proof refers to but does not contain; nothing load-bearing is delegated to it.\par\medskip
\begin{lemma}
\label{lemma-birational-birational}
Let $f : X \to Y$ be a birational morphism of schemes having finitely
many irreducible components over a base scheme $S$. Assume one of the
following conditions is satisfied
\begin{enumerate}
\item $f$ is locally of finite type and $Y$ reduced,
\item $f$ is locally of finite presentation.
\end{enumerate}
Then there exist dense opens $U \subset X$ and $V \subset Y$
such that $f(U) \subset V$ and $f|_U : U \to V$ is an isomorphism.
In particular if $X$ and $Y$ are irreducible, then
$X$ and $Y$ are $S$-birational.
\end{lemma}

\begin{proof}
There is an immediate reduction to the case where $X$ and $Y$ are irreducible
which we omit. Moreover, after shrinking further and we may assume $X$ and
$Y$ are affine, say $X = \Spec(A)$ and $Y = \Spec(B)$.
By assumption $A$, resp.\ $B$ has a unique minimal prime $\mathfrak p$,
resp.\ $\mathfrak q$, the prime $\mathfrak p$ lies over $\mathfrak q$, and
$B_\mathfrak q = A_\mathfrak p$.
By Lemma \ref{lemma-birational-generic-fibres} we have
$B_\mathfrak q = A_\mathfrak q = A_\mathfrak p$.

\medskip\noindent
Suppose $B \to A$ is of finite type, say $A = B[x_1, \ldots, x_n]$.
There exist a $b_i \in B$ and $g_i \in B \setminus \mathfrak q$
such that $b_i/g_i$ maps to the image of $x_i$ in $A_\mathfrak q$.
Hence $b_i  - g_ix_i$ maps to zero in $A_{g_i'}$ for some
$g_i' \in B \setminus \mathfrak q$. Setting $g = \prod g_i g'_i$
we see that $B_g \to A_g$ is surjective. If moreover $Y$ is reduced,
then the map $B_g \to B_\mathfrak q$ is injective and hence
$B_g \to A_g$ is injective as well. This proves case (1).

\medskip\noindent
Proof of (2). By the argument given in the previous paragraph we may
assume that $B \to A$ is surjective. As $f$ is locally of finite presentation
the kernel $J \subset B$ is a finitely generated ideal. Say
$J = (b_1, \ldots, b_r)$. Since $B_\mathfrak q = A_\mathfrak q$
there exist $g_i \in B \setminus \mathfrak q$ such that
$g_i b_i = 0$. Setting $g = \prod g_i$ we see that $B_g \to A_g$
is an isomorphism.
\end{proof}

\end{document}
